\documentclass[pdf]{beamer}

\usepackage{booktabs}
\usepackage{listings}
\usepackage{multicol}
\usepackage{amsmath}
\usepackage{hyperref}

\usepackage{fontawesome}% for Github icon
\usepackage{FiraSans}
\usepackage{FiraMono}


% set theme and colors
\usetheme{ntu}
\setcolor[ntugold]{ntubrick}
\setdarkframecolor{darkthemedeepblue!20!black}
\hypersetup{colorlinks,linkcolor=,urlcolor=\firstcolor}
% Warning : If you set color other than ntu types, the logo in outertheme will disappear. Latex can not change the color of a picture. You might need to work on it yourself (e.g.PS)


% Set up listings env
% basic setups
\lstset{%
    basicstyle=\footnotesize\ttfamily,
    keywordstyle=\color{\secondcolor}\bfseries\ttfamily,
    backgroundcolor = \color{lightgray!20},
    stringstyle = \color{\secondcolor},
    commentstyle=\color{\firstcolor!60},
    numbers=left,
    xleftmargin=1.0ex,
    stepnumber=1,
    numbersep=10pt,
    tabsize=4,
    showspaces=false,
    showstringspaces=false
}


% Set title, author ... etc

\title{NTU \\BEAMER TEMPLATE}
\subtitle{ALL YOU NEED IS SETCOLOR}
\author{Liang-Cheng Chen \& Po-Kang Hsiao}
\date{June 2021}

% include the reference file
\bibliography{reference.bib}% change the path for your .bib
\renewcommand*{\bibfont}{\footnotesize}




\begin{document}
 % if you want to remove the first page logo
 % please go to beamerthementu.sty and comment out first frame of \inserttitlepage command
 \inserttitlepage

\begin{frame}{Outline}

    \tableofcontents
    
\end{frame}


\section{Introduction}

\begin{frame}{Motivation}

    \begin{itemize}
        \item No existing NTU beamer (with logo)
        \item Modern but structural template
        \item Adjustable color
        \item We like dark theme
    \end{itemize}
    
\end{frame}

\begin{frame}{About\, this\, project ...}

    \begin{itemize}
        \item A less beamer-like NTU beamer template 
        \item Easy color switching
        \item Literature importance marking
        \item Colorized description box
    \end{itemize}
    
\end{frame}


\section{Features}

\subsection{Colors}

\begin{frame}{Colors}

\begin{itemize}
    \item Use pre-defined NTU colors
    \item Add new colors to \alert{beamercolorthementu.sty} file
\end{itemize}
\vspace{.5cm}
\begin{mybox}{gray}{This template provides below colors to choose}
    \begin{testcolors}[RGB]
        
        \testcolor{ntubrick}
        \testcolor{ntugold}
        \testcolor{ntudark}
        \testcolor{ntured}
        \testcolor{ntublue}
        
        \testcolor{lakerpurple}
        \testcolor{lakergold}
        
    \end{testcolors}
    
    
    
\end{mybox} 
\scriptsize{*More combination of colors can be found here : \href{https://colorhunt.co}{Color Hunt}}
\end{frame}



\begin{frame}[fragile]{Change Color}

\begin{itemize}

\item Use \alert{setcolor} function to change the main and secondary color of the slides
    
\begin{lstlisting}[style=tex]
% Use NTU theme
\usetheme{ntu}

% Set 2 colors
\setcolor[SecondColor]{MainColor}

% Set main color only
\setcolor{MainColor}
\end{lstlisting}

\item The colors in top bar and footline would be slightly darker
    
\end{itemize}

\end{frame}

\subsection{Itemize/Enumerate}

\begin{frame}{Item and Enumerate}

    \begin{itemize}
        \item The item color 
        \item The subitem would be closer to secondary color
        \begin{itemize}
            \item subitem color 
        \end{itemize}
        
    \end{itemize}
    
    \begin{enumerate}
        \item Same color is set to enumerate as well
        \item The subitem will be closer to secondary
        \begin{enumerate}
            \item subitem color here
        \end{enumerate}
        
    \end{enumerate}
    
\end{frame}

\subsection{Box}

\begin{frame}{Box}

\examplebox{Example box}{
    This \alert{examplebox} command would take 2 arguments:
    \begin{itemize}
        \item box title
        \item box content
    \end{itemize}
    And the box will automatically set to secondary color.
}

\end{frame}

\begin{frame}{Box}
    
\begin{mybox}{\firstcolor}{More flexible box}
Use mybox enviroment for more colors.
\begin{itemize}
    \item Set to firstcolor
    \item Set to arbitary colors
\end{itemize}
\end{mybox}

\begin{mybox}{purple}{Like This box}
mybox\{color\}\{title\}
\end{mybox}
    
\end{frame}

\begin{frame}{More on Box}

\begin{columns}[t]
    
    \begin{column}{0.5\linewidth}
    
        \examplebox{box 1}{This is the box 1.
        \begin{itemize}
            \item First
            \item Second
            \item Third
        \end{itemize}
        
        }
        
        
    \end{column}
    
    \begin{column}{0.33\linewidth}
    
        \examplebox{box 2}{This is the box 2.}
        
    \end{column}

\end{columns}

    
\end{frame}


\begin{frame}{Bayes' theorem}
    \begin{Theorem} %TODO: insert thm number
     \begin{itemize}
        \item Bayes Rule\\
        \[\mathds{P}(A|B) = \frac{\mathds{P}(P|A) \cdot \mathds{P}(A)}{\mathds{P}(B)}\]
        \item Let \(A_1, A_2, \dots , A_K\) be a \textit{partition} of the sample space \(S\). Using the \textit{law of total probability} for \(\mathds{P}(B)\), it then holds that:
        \[
            \mathds{P}(A_j | B) = \frac{\mathds{P}(B|A_j) \cdot \mathds{(A_j)}}{\sum_k \mathds{P}(B|A_k) \cdot \mathds{P}(A_k)}
        \]
    \end{itemize}
    \end{Theorem}
\end{frame}


\begin{frame}{Block}
    
    \begin{block}{Random Block Title}
    \begin{enumerate}

    \item This will be a block \cite{10.1257/aer.20200451}
    \item This will be another citation \cite{10.1257/aer.20190553}
    \item This will also be another citation \cite{Adler:1998:MDL}
    \item This still will be another citation \cite{doi:https://doi.org/10.1002/9780470756461.ch5}
     \end{enumerate}
    \end{block}
    
\end{frame}

\subsection{Overlay}

\begin{frame}{Overlay}

    \begin{itemize}[<+->]
        \item First thing first
        \item Second point 
        \item And if you want to highlight current object
    \end{itemize}
    
    \begin{enumerate}[<+-|alert@+>]
        \item Adding [<|alert@+>] could help coloring
        \item Now we're talking about this point
        \item Test a math expression
        
        \begin{equation*}
            \beta = (X'X)^{-1} \cdot (X'Y)
        \end{equation*}
        
    \end{enumerate}
    
    
\end{frame}

\subsection{Code}

\begin{frame}[fragile]{Code Display}
% now is unable to be flexible adjustable
%TODO : code enviroment

\begin{lstlisting}[basicstyle=\normalsize, language=python]
# Python Code
import torch
import torch.nn as nn

# model architecture
class Classifier(nn.Module):
    def __init__(self):
        super().__init__()
        pass
model = Classifier()

\end{lstlisting}
    
\end{frame}



\subsection{Related Literature}

\begin{frame}{Related Literature}

\examplebox{citation with levels}{
\begin{itemize}
    \item The normal citation would be like this ~\cite{Adler:1998:MDL}
    
    \item If the citation is more realted / important, could set the cite level to A ~\cite[A]{10.1257/aer.20190553}
    
    \item The most important level (level set to "S") will be like this ~\cite[S]{10.1257/aer.20181607}
\end{itemize}

}

\end{frame}



% {
% \setbeamercolor{background canvas}{bg=ntudark} % Bg
% %\setbeamertemplate{frametitle}{\insertframetitle}
% \setbeamercolor{frametitle}{fg = white}
% \setbeamercolor{enumerate}{fg = white}
% \setbeamercolor{itemize}{bg = white}
% \setbeamercolor{itemize/enumerate body}{fg = white}
%  \begin{frame}{Dark Theme}
%  \color{white}
%  \darkthemelogo
%   \begin{itemize}
%         \item he
%         \item boo
%         \begin{block}
%         moomoo
%         \end{block}
%         \item qwe
%     \end{itemize}
     
%  \end{frame}
% }

\begin{darkframe}
\frametitle{Example of a Dark Theme }

    \begin{itemize}
        \item item 1
        \item item 2
        \item test a table
    \end{itemize}
    
    \begin{table}[H]
    \resizebox{.5\textwidth}{!}{%
    \label{tab:tableExample}
    \begin{tabular}{@{}llS@{}}
    \toprule
    \multicolumn{2}{c}{Record} &            \\ \cmidrule(r){1-2}
    User       & Product & Price (\$) \\ \midrule
    UserA      & Product1     & 13.65      \\
               & Product2     & 0.01       \\
    UserB      & Product3     & 92.50      \\
    UserC      & Product4     & 33.33      \\
    UserD      & Product5     & 8.99       \\ \bottomrule
    \end{tabular}%
    }
    \end{table}
    
\end{darkframe}

\begin{darkframe}[summary table]
% \frametitle{Summary Table}
\begin{table}[!htbp] \centering 
  \caption{The iris dataset descriptive statistics.} 
  \label{1} 
\resizebox{\columnwidth}{!}{%
    \begin{tabular}{@{\extracolsep{5pt}}lccccccc} 
    \\[-1.8ex]\hline 
    \hline \\[-1.8ex] 
    Statistic & \multicolumn{1}{c}{N} & \multicolumn{1}{c}{Mean} & \multicolumn{1}{c}{St. Dev.} & \multicolumn{1}{c}{Min} & \multicolumn{1}{c}{Pctl(25)} & \multicolumn{1}{c}{Pctl(75)} & \multicolumn{1}{c}{Max} \\ 
    \hline \\[-1.8ex] 
    Sepal.Length & 150 & 5.843 & 0.828 & 4.300 & 5.100 & 6.400 & 7.900 \\ 
    Sepal.Width & 150 & 3.057 & 0.436 & 2.000 & 2.800 & 3.300 & 4.400 \\ 
    Petal.Length & 150 & 3.758 & 1.765 & 1.000 & 1.600 & 5.100 & 6.900 \\ 
    Petal.Width & 150 & 1.199 & 0.762 & 0.100 & 0.300 & 1.800 & 2.500 \\ 
    \hline \\[-1.8ex] 
    \end{tabular}%
}
\end{table} 
\end{darkframe}




\begin{darkframe}
\frametitle{Graph is more vivid within the dark theme}
\centering{\includegraphics[width=9cm,keepaspectratio]{logos/cat.jpeg}}
% The pic is from unsplash where one will not encounter copyright problems.
\end{darkframe}

\begin{frame}[t, allowframebreaks]{Reference}

\printbibliography
    
    
\end{frame}

\section{Project}


\begin{frame}{Download}

    To use this template, you could download the whole project from
    the overleaf or github page.\\
    % \url{https://github.com/liangchenn/NTU-Beamer}\\
    \href{https://www.overleaf.com/read/qppxwpysgjmv}{\beamergotobutton{overleaf link}}~~
    \href{https://github.com/liangchenn/ntu-beamer-template}{\faGithub}
    
\end{frame}


\begin{frame}{Still working on...}

\begin{itemize}
    \item Individually stylize the references
    \begin{itemize}
        \item We can not change the font by setbeamerfont and we do not know why 
        \item We expect to trasfer our newcommand cite's second argument into the coloring or font size or ordering in printing reference. 
    \end{itemize}
    \item Progress bar on headline.
    \item Copyright of the icon NTU, and reference from other projects we should refer to.
    
\end{itemize}
    
\end{frame}


\end{document}
